\section{Banked Cloze}
\subsection{Unit 1}

\begin{question}{Banked Cloze}

Cities are getting smarter every day, so are their transit options. One emerging option is autonomous vehicles (AVs), or simply cars with self-driving  \textbf{\textcolor{red}{capability}}. However, as experts say, before we enter a reality where AVs are  \textbf{\textcolor{red}{embedded}} into our daily lives, we will see the continued rise of electric vehicles (EVs) all over the world.

Besides saving costs for consumers, EVs can bring many \textbf{\textcolor{red}{sustainability}} benefits to the environment because they use electricity instead of gas to drive their wheels. Many large cities are trying to encourage the use of EVs by installing public charging stations throughout the city and \textbf{\textcolor{red}{streamlining}} the vehicle registration process. At the same time, many automakers are working to \textbf{\textcolor{red}{enhance}} the charging speed of EVs and to make EVs more \textbf{\textcolor{red}{affordable}}.

As the use of EVs rises, greater \textbf{\textcolor{red}{demand}} will be placed on the power industry. In some areas where energy resources are stressed already, the increasing use of EVs could well \textbf{\textcolor{red}{disrupt}} the electrical supply system.

EVs can also impact urban planning. Thus, consideration of commute options, location of charging stations, and open data should be part of a city's future mobility plans. It is important for automakers to partner with urban planners and utility companies to provide \textbf{\textcolor{red}{optimal}} experiences for consumers and residents. Imagine one day a driverless EV picks you up and drops you off, and then \textbf{\textcolor{red}{proceeds}} to the location of the next passenger without ever stopping to recharge because of wireless power transfer from nearby energy sources. The technology is on the way, but we need to wait for our cities and energy industry to get ready as well.

\end{question}

\newpage

\subsection{Unit 2}

\begin{question}{Banked Cloze}

When I was young, I had a very narrow understanding of beauty. I used to \textbf{\color{red}associate} it with only physical things. It was about what the eye could see: soft hair, clear skin,  \textbf{\color{red}stylish} clothes, a good figure ... I think some other people may also have gone through this phase of wanting the perfect shell defined solely in terms of  \textbf{\color{red}physical} standards.

When I first entered the modeling industry, I was always  \textbf{\color{red}striving} to conform to others' criteria of beauty. I would think, ``Well, if they say I'm beautiful, then I am. But if they don't, then I'm not.'' Even today, I can't believe how  \textbf{\color{red}ridiculous} my thoughts were.

I want to correct this narrow idea that our beauty comes from the notions and approval of other people. When we allow others to determine what is beautiful, we fail to express and accept our individuality, our unique physical and personal \textbf{\color{red}traits}. I think we should ask ourselves the question ``What is beautiful to me?'' and construct our own \textbf{\color{red}concepts} of beauty that are more comprehensive and more meaningful.

As time goes on, I have gradually expanded my definitions of beauty, without limiting them to what lies on the outside. I have started to realize that physical beauty won't last forever because we humans are all  \textbf{\color{red}mortal} and will grow old inevitably. I have met many strong and courageous women, for example, a friend who \textbf{\color{red}survived} breast cancer and a cousin who works hard to provide for her child. I have also met many women who are devoted to their careers. Their confidence radiates from their \textbf{\color{red}dedication} to their professional pursuits. These women are examples of beautiful actions and thoughts. They embrace being themselves.

\end{question}

\newpage

\subsection{Unit 3}

\begin{question}{Banked Cloze}

Welcome to our special edition "Building Better Businesses". As is well recognized, the business world is changing. For the new generation of entrepreneurs, setting up a Silicon Valley social media venture isn't their dream. Instead, their \textbf{\color{red}priority} is clear: building a company that has a purpose at its heart. These founders speak with \textbf{\color{red}clarity} that they want their companies to have a social or environmental impact as well as to make money. They want to play a role in their communities, \textbf{\color{red}strike} a balance between work and private life, and invest for the long term rather than for quick profits.

At the same time, older people are becoming more active in entrepreneurship, too, as they find they still \textbf{\color{red}harbor} the dream of setting up their own business and now have the confidence to take the risk. With people living longer and starting to seek more \textbf{\color{red}accomplishments} in later life, this trend will surely continue to grow in the future. Plus, it is noted that older founders are more \textbf{\color{red}motivated} to succeed in business than their younger counterparts (相对应者) because they believe they have more experience and wisdom.

Our special edition intends to help you build a better business, whether you are an ambitious individual who \textbf{\color{red}aspires} to launch an independent company and strives to make a(n) \textbf{\color{red}positive} impact on the world, or whether you are already a business owner and are looking for inspiration to \textbf{\color{red}diversify} your services. Whatever your story, we hope you will find all the \textbf{\color{red}ingenious} ideas, helpful tips, instructions, and insights you need in this guide. So, why not begin your exploration with us now?
	
\end{question}

\newpage

\subsection{Unit 4}

\begin{question}{Banked Cloze}

Farming invariably interferes with the habitats of plants and animals. However, this does not necessarily mean that agriculture will put biodiversity in \textbf{\color{red}{jeopardy}} and that the two can never come together. In fact, quite the opposite is true. Sustainable cultivation of agricultural land for food and animal feed enables us to \textbf{\color{red}{conserve}} biodiversity.

Compared with the massive population growth in the last few decades, the expansion of arable (可耕的) land over the same period has been small. And there are limits to further expansion. A large \textbf{\color{red}{proportion}} of the earth's surface like deserts is not suitable for crop cultivation, and other areas are \textbf{\color{red}{utilized}} by humans for roads and buildings. Some land that is rich in biodiversity needs to be protected and thus should not be \textbf{\color{red}{converted}} into arable land. The tropical rain forests, for example, have the highest species diversity in the world, and changing this land for crop cultivation would be \textbf{\color{red}{devastating}} to the habitats and, indeed, the existence of the species in these places.

By 2050, global demand for food will have increased substantially. This is one of the \textbf{\color{red}{monumental}} challenges facing agriculture today: How do we balance the rising demand for food and the need to maintain biological \textbf{\color{red}{integrity}} and diversity, now and in the future?

Efficient and effective use of land will be \textbf{\color{red}{imperative}} to the increase of food production and preservation of natural animal and plant habitats. Achieving this will depend, to a large extent, on the use of modern and \textbf{\color{red}{innovative}} agricultural methods. If these methods are successfully developed, agriculture and biodiversity will be able to coexist in greater harmony.

\end{question}

\newpage

\subsection{Unit 5}

\begin{question}{Banked Cloze}

Jack MacDonald, a Seattle man, left a fortune worth 
\$187.6 million to charity after living a frugal and humble life. The {\color{red}\textbf{donation}} was the largest charitable gift in Washington State in 2013.

MacDonald struck people as a man without much money. He {\color{red}\textbf{clipped}} coupons to save money on groceries, wore sweaters with holes, and took a bus to attend alumni luncheons. He once bought so many cans of frozen juice at a bargain price that he had to {\color{red}\textbf{purchase}} a new freezer to hold them.

MacDonald left 40 percent of his fortune to Seattle Children's Research Institute to support a variety of {\color{red}\textbf{innovative}} research in childhood diseases, although he had no {\color{red}\textbf{offspring}} of his own. He was born in British Columbia, Canada, grew up in Seattle, and worked as a lawyer for three decades. He did not get married until he was in his 50s.

MacDonald made a huge fortune by investing in the stock market. After his parents died, he received a large family {\color{red}\textbf{inheritance}}. He never thought about buying any luxuries although he had {\color{red}\textbf{enormous}} wealth at his disposal. Material possessions weren't important to him.

MacDonald might have been eccentric in some ways, but he always acted on his {\color{red}\textbf{convictions}} to do good with his wealth. Over the years, he contributed to hundreds of causes and organizations. He also sent anonymous gifts to a small town in Canada where his grandparents lived after they {\color{red}\textbf{immigrated}} there from Scotland.

In July 2013, MacDonald was hospitalized due to a serious head injury that would eventually lead to his death. In the two months before he died, he insisted on using inexpensive ordinary drugs so that his treatment wouldn't {\color{red}\textbf{incur}} additional costs. He once said he would like to be remembered as a philanthropist (慈善家). With his generosity and selflessness, surely he has been.
	
\end{question}

\newpage

\subsection{Unit 6}

\begin{question}{Banked Cloze}

It's true that greenhouses consume resources, but making the correct choices when you design, build, and run your greenhouses can make your growing operation more sustainable.

Typical greenhouse operation consumes {\color{red}\textbf{lavish}} amounts of plastic, from the plastic film covering on many greenhouses, to the plastic packaging of soil and fertilizers. Choosing recyclable materials, both in the {\color{red}\textbf{construction}} of greenhouses and in their everyday operation, is one effective way to run sustainable greenhouses.

By their very nature, greenhouses are good at collecting energy from the {\color{red}\textbf{glittering}} sun and using it to warm or heat their {\color{red}\textbf{interiors}} to the appropriate temperature. Passive solar design, which doesn't use mechanical or electrical {\color{red}\textbf{apparatuses}} but simply a building's own structure to capture and store solar energy, is usually most people's choice for their greenhouses. Meanwhile, allowing for adequate control of {\color{red}\textbf{ventilation}} when you design greenhouses minimizes the need for electric fans and other energy-consuming cooling solutions. A good greenhouse design also considers the installation of thermal-banking system, a process that uses the sun to heat water and then {\color{red}\textbf{stores}} the water underground, where it can be used to heat the greenhouse in cooler weather. Though the thermal-banking project sounds a little {\color{red}\textbf{intricate}} due to the extra design and work it requires, it helps make the operation of the greenhouse more sustainable.

The operation of a greenhouse consumes energy and uses a lot of water as well. To {\color{red}\textbf{resolve}} this problem, you can replace traditional lighting with more energy-efficient LED lighting, and build a rainwater collection and storage system. These methods can significantly {\color{red}\textbf{slash}} energy and water costs though some more investments are needed at first.
	
\end{question}

\newpage